\section{Antecedentes}
%% Historia
%% Auge en la vida actual
%% Crisis de confianza, XAI 
%% Amarre a tema de tesis

La inteligencia artificial (IA) es una disciplina científica cuyo origen conceptual se entrelaza con la evolución de las
matemáticas que la sustentan. Desde una perspectiva teórica, uno de los fundamentos más remotos se atribuye a Gottfried Wilhelm Leibniz, 
quien en 1676 desarrolló la regla de la cadena; esta fórmula del cálculo diferencial constituye la base del algoritmo de propagación hacia atrás 
(\textit{backpropagation}) empleado en el aprendizaje de las redes neuronales modernas, además de anticipar los principios de las ciencias
de la computación. Por otra parte, su origen operativo se vincula al trabajo de Adrien-Marie Legendre, quien en 1805 publicó 
el método de los mínimos cuadrados (conocido actualmente como regresión lineal). Este modelo matemático es considerado el precursor estructural 
directo de las neuronas artificiales y las redes de dos capas contemporáneas \parencite{schmidhuber2022annotated}.
\medskip

Asimismo, el interés por la IA dentro de la cultura popular ha sido impulsado históricamente por narrativas sobre máquinas capaces de emular el razonamiento humano. 
Un ejemplo célebre es el Turco Mecánico de 1770, un autómata que presuntamente jugaba al ajedrez de forma autónoma e incluso venció a figuras como Napoleón Bonaparte 
y Benjamin Franklin, aun cuando en realidad era operado por un maestro de ajedrez oculto en su interior. En la literatura, las obras de Isaac Asimov representan el referente 
más influyente, explorando las relaciones éticas y operativas entre humanos y máquinas a través de las tres leyes de la robótica, las cuales han inspirado a generaciones de 
científicos en los campos de la robótica, la IA y las ciencias de la computación \parencite{haenlein2019brief}.
\medskip

En el ámbito científico, Alan Turing destaca como una de las figuras más influyentes; su trabajo en la decodificación de la máquina Enigma estimuló su interés por la cognición mecanizada. 
En 1950, Turing propuso el juego de la imitación, popularmente conocido como la prueba de Turing, un criterio diseñado para evaluar la inteligencia de una máquina con base en la 
incapacidad de un evaluador humano para distinguir si interactúa con otra persona o con un sistema computacional \parencite{turing1950computing}. Otro hito fundacional fue el proyecto de investigación de verano de 
Dartmouth sobre Inteligencia Artificial (\textit{Dartmouth Summer Research Project on Artificial Intelligence}, DSRPAI) en 1956. En este encuentro se acuñó formalmente el término 
«inteligencia artificial» y se reunieron los investigadores que definirían las bases metodológicas de esta nueva área del conocimiento.
\medskip

La inteligencia artificial floreció en las décadas subsecuentes impulsada por los avances tecnológicos y 
computacionales, los cuales permitieron el desarrollo de algoritmos con capacidades de abstracción y procesamiento 
de información cada vez más complejas. En la actualidad, esta tecnología es más accesible que nunca y se encuentra 
integrada en múltiples aspectos de la vida cotidiana y laboral. Por ejemplo, los automóviles autónomos capaces de 
operar sin conductor ya constituyen una realidad comercial, con servicios de robotaxis que completan cientos de miles 
de viajes semanales. Asimismo, en el ámbito de la medicina, la inteligencia artificial ha acelerado la investigación 
clínica y ha demostrado capacidades operativas que equiparan o superan el rendimiento humano en tareas diagnósticas 
especializadas, tales como la detección temprana de cáncer. 
\medskip

Por otra parte, dentro del entorno corporativo, un número creciente de organizaciones ha adoptado herramientas de IA, 
registrando incrementos significativos en la productividad y logrando mitigar la brecha de competencias entre 
trabajadores de distintas cualificaciones operativas. Se proyecta que esta adopción se intensifique en los próximos 
años, a medida que la infraestructura computacional requerida optimice su eficiencia y reduzca sus costos, y el 
creciente interés global continúe impulsando la investigación y el desarrollo hacia modelos arquitectónicos más 
robustos y avanzados.
\medskip

A pesar de estos avances, la adopción masiva de la inteligencia artificial no está exenta de controversias y reservas 
estructurales. El impacto ambiental, medido a través de las emisiones de carbono derivadas del entrenamiento de modelos 
complejos, se ha incrementado significativamente en los últimos años, al igual que el número de incidentes operativos 
y éticos reportados en la industria. En consecuencia, la opinión pública permanece dividida respecto a la confiabilidad 
de estos sistemas; gran parte del escepticismo se centra en el comportamiento ético de las corporaciones tecnológicas, 
lo que ha impulsado a un número creciente de naciones a promulgar marcos regulatorios y legislaciones específicas para 
el control de la IA.
\medskip

Dentro de este panorama, una de las mayores preocupaciones reside en el desarrollo y despliegue de una Inteligencia 
Artificial Responsable (\textit{Responsible AI} o RAI). Los criterios estandarizados para evaluar la RAI siguen siendo 
limitados, y las estrategias para mitigar riesgos críticos —tales como la falta de precisión, el incumplimiento 
normativo y las vulnerabilidades en ciberseguridad— aún carecen de una adopción generalizada. Asimismo, las métricas 
diseñadas para auditar la veracidad de las salidas no se han consolidado en los flujos de producción. Como resultado, 
múltiples modelos vigentes continúan exhibiendo sesgos de imparcialidad que afectan desproporcionadamente a minorías 
étnicas y mujeres, al tiempo que persiste una profunda incertidumbre sobre el impacto sistémico del uso de la IA en 
la diseminación de información falsa durante procesos electorales a nivel global \parencite{maslej2025artificial,bughin2017artificial}.
\medskip

